\section{Three-root transport: eleven positive variables, degree 325}
\label{rt11:section}

There is a division-free polynomial in eleven strictly positive integer
unknowns of exact total degree $325$, evaluated by $113$ binary arithmetic
operations, whose lack of positive solutions is equivalent in ZFC to
$\operatorname{Con}(\mathrm{ZFC})$. At the pinned repository state the
eleven-positive equation in package~52 has degree $954$ and the same
$113$-operation cost. The construction below removes $629$ degrees while
retaining that operation count. Its two residuals have degrees $312$ and
$13$ and together require $110$ operations.

The new ingredient is a particular combination of the normalized positive
auxiliary reconstruction, the weighted first carrier, and an elementary
denominator argument. Conjugate norms, positive majorants, radical
separation, quotient absorption and unit-product finalizers already occur
in package~52 and its predecessors; those general methods are not claimed
as new. In package~52 the auxiliary radicand depends on the signed strong
Pell root and the main Pell root is eliminated inside the same norm.
Here the normalized auxiliary expression depends only on its square.
The first, strong and auxiliary roots therefore enter a direct octic
sum norm, while the inexpensive adaptive main equation remains separate.

\subsection{Fixed data and the literal equation}

Use precisely the coherent tuple $(\beta,K_0,M_C,M_F)$ of
Section~\ref{sec:root-transfer}, with $\beta=2^{590069}-1$ and the already
shifted $M_F$. These are fixed integer numerals, not parameters.
Their payloads and source hashes are supplied again in this component.
The authentic mask inequalities used below are
\[
 0<M_C<\beta,\qquad 0<M_F<2\beta,\qquad K_0>0.
\]
The eleven unknowns are exactly
\[
 j,F,\alpha,i,T,X,Y,n,\eta,\zeta,\gamma\in\mathbb Z_{>0}.
\]
Define polynomial registers
\begin{align*}
 J&=3j+1,&q&=\beta J+1,&C&=q-F-\alpha,\\
 E&=XY,&k&=\eta+\zeta,&c&=kY+\eta,\\
 a&=E+Y,&H&=4a+3,&\Delta&=a^2+H=(a+2)^2-1,\\
 \ell&=EkY,&M&=q^3(q-1),&B_t&=(qK_0+X)C-qF,\\
 D_\tau&=\ell(\ell+k)+1.&&&
\end{align*}
Set
\begin{align}
 R&=\bigl((q-1)(q-F)+\alpha\bigr)(q^2-1)+(M_C+qM_F)J,
       \label{rt11:R}\\
 \widehat R&=q^3(q+F+\alpha),\nonumber
\end{align}
and then
\begin{align}
 I&=iq(q-1)c^2,&D_f&=1+\Delta I^2,&Q&=\Delta(D_f-1),
       \label{rt11:strong}\\
 \delta&=Q-1,&V&=D_f(cT-R)-c,&D_y^{\rm num}&=QV^2-1,
       \nonumber\\
 \overline V&=D_f(cT+\widehat R),&L&=3\overline V,&W&=1+(q-1)L.
       \label{rt11:majorant}
\end{align}
The first transport register is
\begin{equation}
 \overline Z=M(\ell+k+n+L)+(q-1)Y-q^3WB_t.
 \label{rt11:Z}
\end{equation}
The adaptive main carrier is unchanged:
\begin{align}
 r_0&=k-R-1,& S&=4qY,& U&=qS,\nonumber\\
 D_*&=U(X+ac)+H(U\gamma+SX-r_0),\nonumber\\
 A_1&=D_*^2-U^2(1+\Delta c^2).
 \label{rt11:main}
\end{align}

The following schedule evaluates the octic norm using only integer
polynomial operations:
\begin{align}
 \mathcal A&=M^2D_\tau,&
 \mathcal B&=M^2D_f,&
 \mathcal C&=M^2D_y^{\rm num},\nonumber\\
 T_n&=\delta^2\mathcal A\mathcal B,&
 R_n&=\delta\mathcal C,&
 J_n&=\delta(\overline Z^2-\mathcal A-\mathcal B)-\mathcal C,
       \nonumber\\
 N_*&=J_n^2-4\bigl(T_n+R_n(\mathcal A+\mathcal B)\bigr),
       \nonumber\\
 N_8&=N_*^2-64T_nR_n\overline Z^2.
 \label{rt11:norm}
\end{align}
The equation is
\begin{equation}
 \boxed{\ P_{11}=N_8+A_1+N_8A_1=0.\ }
 \label{rt11:equation}
\end{equation}
All the displayed registers are abbreviations, with no additional
quantified witnesses. The supplied JSON is a literal binary
$+,-,\times$ circuit; even multiplication by a small integer is charged.
The rational quantity
\[
 D_y=\frac{D_y^{\rm num}}\delta
\]
is used only in the proof and is never evaluated by the polynomial circuit.

\subsection{An elementary rational-root extraction lemma}

\begin{lemma}
\label{rt11:rational}
Let $d_1,d_2,d_3$ be distinct positive rational numbers. Suppose signed
real roots $x_s^2=d_s$ have a nonzero rational sum $t=x_1+x_2+x_3$.
Then all three $x_s$ are rational.
\end{lemma}
\begin{proof}
Put
\[
 u=\frac{t^2-d_1-d_2-d_3}{2},\qquad
 v=\frac{u^2-d_1d_2-d_1d_3-d_2d_3}{2t}.
\]
These are rational. Expansion gives
$u=x_1x_2+x_1x_3+x_2x_3$ and $v=x_1x_2x_3$. Consequently
\[
 d_1t+v=x_1(x_1+x_2)(x_1+x_3),\qquad
 x_1=\frac{d_1(d_1+u)}{d_1t+v}.
\]
The denominator is nonzero. Indeed, $x_1\ne0$, while either
$x_1=-x_2$ or $x_1=-x_3$ would contradict distinctness of the radicands.
The two cyclic formulas recover $x_2,x_3$ as rational numbers as well.
This proof accommodates every coincidence or dependence among square
classes; it assumes no independent specialized Galois automorphisms.
\end{proof}

\subsection{Unconditional radical separation and the positive guard}

All bounds in this subsection hold at every positive input tuple, before
assuming any residual vanishes. In particular they do not assume $C>0$,
$q^3\mid Y$, or a recovered Pell index. We have
\[
 q\ge16,\quad k\ge2,\quad c\ge3,\quad a\ge2,\quad
 \Delta\ge15,\quad I\ge c^2,\quad \ell<ac,\quad k<c.
\]
Thus
\[
 \ell(\ell+k)<a(a+1)c^2<\Delta I^2,
 \qquad 1<D_\tau<D_f<Q.
\]
Write $t=\sqrt{D_\tau}$ and $f=\sqrt{D_f}$ for positive real roots.
Then
\begin{equation}
 t<\ell+k<(a+1)c,
 \qquad f>(a+1)I\ge(a+1)c^2\ge3c^2.
 \label{rt11:tf}
\end{equation}

The integer $b=cT-R$ determines two cases. If $b=0$, then $V=-c$.
If $b\ne0$, then $|V|\ge D_f-c>1$. Hence $|V|>1$ in both cases and
\begin{equation}
 D_y=V^2+\frac{V^2-1}{Q-1}>V^2>1.
 \label{rt11:Dy}
\end{equation}
The radicands $D_\tau,D_f,D_y$ are pairwise distinct. For if
$D_y=D$ with $D\in\{D_\tau,D_f\}$, then
\[
 Q(D-V^2)=D-1,
\]
which is impossible because $0<D-1<Q$.

Let $y=\sqrt{D_y}>0$. If $b=0$, then $Q>c^2$ implies
$D_y<c^2+1$ and $y<c+1$, whereas \eqref{rt11:tf} gives
\[
 f-t>(a+1)c(c-1)>c+1.
\]
Thus $f>t+y$. If $b\ne0$, the bound $f>3c^2>c+2$ implies
$f^2-c>2f$, so
\[
 y>|V|\ge D_f-c>2f>f+t.
\]
In either case one root exceeds the sum of the other two. Therefore no
signed sum $\pm t\pm f\pm y$ is zero.

It remains to check that the polynomial $L$ bounds the two roots to be
transferred out of the positive coordinate. The authentic masks give
\[
 (M_C+qM_F)J<(q-1)(1+2q)<2q^2.
\]
The triangle inequality applied to \eqref{rt11:R} then gives
\[
 |R|<q^4+Fq^3+(\alpha+2)q^2
     <q^3(q+F+\alpha)=\widehat R,
\]
because $\alpha(q-1)>2$. Both sides are integers, and $D_f>c$; hence
\[
 |V|\le D_f(cT+|R|)+c
       <D_f(cT+\widehat R)=\overline V.
\]
Also $\overline V>D_f>f$ and, by \eqref{rt11:Dy},
$y<2|V|<2\overline V$. We obtain the uniform bounds
\begin{equation}
 L=3\overline V>f+y,
 \qquad L>c>k.
 \label{rt11:Lbound}
\end{equation}

\subsection{What the cleared norm forces}

The universal octic identity is
\begin{align}
 F_8(x;a,b,c)
 &:=\bigl((x^2-a-b-c)^2-4(ab+ac+bc)\bigr)^2-64abcx^2
       \nonumber\\
 &=\prod_{\epsilon_1,\epsilon_2,\epsilon_3\in\{-1,1\}}
       \bigl(x-\epsilon_1\sqrt a-\epsilon_2\sqrt b
                    -\epsilon_3\sqrt c\bigr).
 \label{rt11:genericnorm}
\end{align}
This is a polynomial identity, verified by multiplying the eight formal
conjugates. The center in \eqref{rt11:norm} also equals
\[
 \delta^2\bigl((\overline Z^2-\mathcal A-\mathcal B)^2
                       -4\mathcal A\mathcal B\bigr)
 -2\delta\mathcal C(\overline Z^2+\mathcal A+\mathcal B)
 +\mathcal C^2.
\]
Consequently
\begin{equation}
 N_8=\delta^4F_8
       (\overline Z;M^2D_\tau,M^2D_f,M^2D_y).
 \label{rt11:clearing}
\end{equation}
All three specialized radicands are positive, and $M,\delta>0$.
If $N_8=0$, one of its eight real factors vanishes. There are signed
roots satisfying
\begin{equation}
 \frac{\overline Z}{M}=\tau_*+f_*+y_*,
 \qquad \tau_*^2=D_\tau,\quad f_*^2=D_f,\quad y_*^2=D_y.
 \label{rt11:signedsum}
\end{equation}
The preceding separation proves that this rational sum is nonzero and
its three radicands are distinct. Lemma~\ref{rt11:rational} makes all
three signed roots rational. Since $D_\tau,D_f$ are integers,
rational-square integrality yields $\tau_*,f_*\in\mathbb Z$.

The auxiliary root becomes integral for a different, elementary reason.
Let $d$ be its positive reduced denominator. Equation
\eqref{rt11:signedsum} gives $d\mid M$, and
$y_*^2=D_y^{\rm num}/\delta$ gives $d^2\mid\delta$.
Every prime divisor of $M=q^3(q-1)$ divides $q(q-1)$ and therefore divides
$I=iq(q-1)c^2$ and $Q=\Delta^2I^2$. It cannot divide $\delta=Q-1$.
Thus
\[
 \gcd(M,\delta)=1,\qquad d=1,
\]
and $y_*$ is an integer. This uses only the prime support of $M$; no
claim that $M\mid Q$ is needed.

\subsection{Soundness and completeness}

\begin{theorem}
\label{rt11:theorem}
The equation \eqref{rt11:equation} has a strictly positive solution if
and only if the inherited normalized compiler has a strictly positive
solution. Hence, using the same compiler and refutation calibration,
ZFC proves that \eqref{rt11:equation} has no positive solution if and
only if $\operatorname{Con}(\mathrm{ZFC})$.
\end{theorem}
\begin{proof}
First, the single equation is equivalent to $N_8=A_1=0$. Indeed,
\[
 P_{11}=(1+N_8)(1+A_1)-1.
\]
Both residuals are squares modulo $4$: this follows directly from
\eqref{rt11:norm} and from $4\mid U$ in \eqref{rt11:main}.
At a zero of $P_{11}$, the two integer factors are therefore both $1$;
the alternative of both being $-1$ would force a residual to be
$-2\equiv2\pmod4$. The converse is immediate.

For soundness, recover the signed integral roots in
\eqref{rt11:signedsum}, and set
\[
 i_{\rm old}=iq(q-1),\quad f_{\rm old}=|f_*|,\quad
 y_{\rm old}=|y_*|,\quad u=y_{\rm old}-V,\quad
 r=n+L-f_*-y_*.
\]
Equation \eqref{rt11:Dy} gives $y_{\rm old}>|V|$, so $u>0$.
The uniform bound \eqref{rt11:Lbound} gives $r>0$ for every sign choice.
The strong equation is
\[
 f_{\rm old}^2-\Delta(i_{\rm old}c^2)^2=1.
\]
The auxiliary root equation is exactly
\[
 y_{\rm old}^2-Q(y_{\rm old}^2-V^2)=1,
\]
which is the normalized auxiliary residual with $y_{\rm old}=V+u$.
Finally the first transport register with weight $h_*=L$ becomes
\[
 Z=M(\ell+k+r)+(q-1)Y-q^3WB_t
   =\overline Z-M(f_*+y_*)=M\tau_*.
\]
It satisfies $Z^2=M^2D_\tau$. The weighted transport and adaptive main
carrier theorem from Section~\ref{sec:root-transfer} now applies, with
all its witnesses positive, and recovers a normalized compiler solution.
The signed roots must be retained in the definition of $r$; absolute
values are taken only in the strong and auxiliary witnesses.

For completeness, start with a canonical normalized compiler solution.
Keep its outer quantities, including $q,X,Y,k,c,a,R$ and its first
positive Pell root $\tau$. The canonical facts needed from the weighted
transport theorem are
\[
 R\equiv3\pmod4,\quad c=\psi_{a+2}(R),\quad
 t_0=\frac{B_t}{q-1}=q(z_{\rm source}-1)\ge q,
 \quad s=\frac{Y}{q^3}<k,\quad \tau>\ell.
\]
We must choose the strong Pell solution with the additional divisibility
$q(q-1)c^2\mid I$; this is not an assumption about a preselected source
witness. Put $m_0=q(q-1)c^2$ and consider
\[
 U_0=\begin{pmatrix}a+2&\Delta\\1&a+2\end{pmatrix}.
\]
Its determinant is $1$, so its reduction belongs to the finite group
$\mathrm{SL}_2(\mathbb Z/m_0\mathbb Z)$. A positive power is the
identity. Since the top-left and lower-left entries of $U_0^v$ are
$\chi_{a+2}(v)$ and $\psi_{a+2}(v)$, a positive multiple of that order
gives positive integers $f,I$ with
\[
 f^2-\Delta I^2=1,\qquad m_0\mid I.
\]
Set $i=I/m_0>0$. The positive auxiliary reconstruction lemma in the
pinned package~61 report (its label \texttt{lem:auxcomplete}) applies to
every such strong solution: its hypotheses are $R\equiv3\pmod4$,
$c=\psi_{a+2}(R)$ and $c\mid I$. It supplies positive $T,u$ and
$y=V+u>0$ satisfying the normalized auxiliary equation. Thus the stronger
divisibility preserves the required reconstruction promises.

Only after choosing these roots and $T$ do we compute
$\widehat R,\overline V,L,W$. Choose
\[
 r=Wt_0+\tau-\ell-k-s.
\]
Then
\[
 r>qW-2k=q+q(q-1)L-2k>L,
\]
using $q\ge16$ and $L>k$. Consequently
\[
 n=r-L+f+y>0.
\]
It follows that $\overline Z=M(\tau+f+y)$, so the norm vanishes.
The unchanged adaptive main construction supplies positive $\gamma$
with $A_1=0$. Its canonical repunit choice also gives $J=3j+1$ with
$j>0$, as in Section~\ref{sec:root-transfer}. This completes the lift.
The order of choices was the strong pair, the auxiliary pair, the
majorant, and then the transport coordinate; there is no circular
choice of witnesses or unproved enlargement of a special index.
\end{proof}

\subsection{Exact degree, circuit cost and verification boundary}

The ordinary total degrees in the eleven supplied unknowns are
\[
\begin{array}{c|rrrrrrrrr}
\text{register}&I&D_f&Q&V&L&W&\overline Z&D_y^{\rm num}&\mathcal C\\\hline
\text{degree}&7&18&22&22&22&23&28&66&74
\end{array}
\]
and $\deg\mathcal A=16$, $\deg\mathcal B=26$.
In the expanded expression for $N_*$, the term
$\delta^2\overline Z^4$ has degree $156$, while the other two groups
have degrees at most $152$ and $148$. The norm tail has degree $238$,
below $2\cdot156$. Thus the top part of $N_8$ is
$Q_0^4\overline Z_0^8$, of degree $312$.

For an explicit certificate, set
\[
 q_0=3\beta j,\qquad k_0=\eta+\zeta,\qquad
 B_{t,0}=(K_0q_0+X)(q_0-F-\alpha)-q_0F.
\]
The relevant leading forms are
\begin{align*}
 Q_0&=i^2q_0^4k_0^4X^4Y^8,\\
 \overline Z_0&=-3i^2q_0^{11}k_0^4X^2Y^6
                      (q_0+F+\alpha)B_{t,0},\\
 (A_1)_{13}&=32q_0^2k_0X^2Y^4
                   \bigl(4q_0^2Y\gamma+q_0^3(q_0-F)\bigr).
\end{align*}
The apparent degree-$14$ terms of $A_1$ cancel; the last expression is
verified by exact sparse polynomial expansion, not by gate-degree
propagation. Therefore the degree-$325$ part of $P_{11}$ is
\begin{align*}
 2^5 3^8\,i^{24}q_0^{106}k_0^{49}X^{34}Y^{84}
 (q_0+F+\alpha)^8B_{t,0}^8
 \bigl(4q_0^2Y\gamma+q_0^3(q_0-F)\bigr).
\end{align*}
Its coefficient of
\begin{equation}
 j^{108}i^{24}\eta^{49}X^{42}Y^{85}\alpha^{16}\gamma
 \label{rt11:coefficientmonomial}
\end{equation}
is exactly
\[
 2^7 3^8(3\beta)^{108}>0.
\]
Indeed the $\gamma$ term fixes the final factor, and minimum $j$ degree
reduces the two remaining parenthesized powers to
$X^8(F+\alpha)^{16}$, whose $\alpha^{16}$ coefficient is $1$.
This certifies degree $325$ on the actual fixed coefficient slice; no
specialization of $K_0,M_C,M_F$ can cancel the selected coefficient.

The literal circuit has $67$ multiplications, $29$ additions and
$17$ subtractions, for $113$ operations. Removing the three-operation
finalizer leaves $110$ operations for the two-residual system. The
normal and optimized-Python replay receipts check byte-identical circuit
regeneration, input and gate liveness, every polynomial register contract,
the formal octic identity, denominator clearing, rational-extraction
identities, the main-Pell cancellation, the exact degree certificate,
and the complete fixed payloads and retained source hashes.
An independent audit uses a separate sparse-ring implementation and
does not import the construction builder.

These checks certify finite algebra and data integrity. The existence
reduction is the proof above together with the preceding weighted carrier
theorem; the unbounded compiler and K215 refutation calibration remain
inherited mathematical dependencies. No end-to-end proof-assistant
certificate, global minimality theorem, or proof of consistency is claimed.
